\section{Uniqueness of the opposing region}
\label{sec:area_law_unique_side_opponents}

The preceding existence assertion can be sharpened to uniqueness of the
opposing fine cell or dummy neighborhood. This identifies the region
containing an elementary side, before assigning primary or fan-run labels;
see \cite[Section~11]{OpenAI2026AreaLaw}.

% Source: 10-geometry.tex, lines 299--310 for elementary geometry and
% lines 299--316 for the opposing region; prop:two-families,
% revision adc7f1241b42e322a6451854ab7e4b4c146bf78a.
% These are scoped assertions about actual cells, not the full proposition.

\begin{theorem}[Elementary sides lie on a nondegenerate square boundary]
    \label{thm:area_law_elementary_side_boundary_geometry}
    \lean{TNLean.PEPS.AreaLaw.Geometry.cellFan_elementary_geometry}
    \leanok
    \uses{def:area_law_cell_fan,def:area_law_cell_corner}
    Fix an origin $o\in\mathbb R^2$, a natural exponent $\ell$, a signed
    cell index $z\in\mathbb Z^2$, and any optional midpoint subdivision
    of the sides of $Q_\ell(z)$. Put $t=2^\ell$, $p=o+tz$ and
    $q=p+(t,t)$. For the ordered endpoints $a,b$ of every elementary
    segment, one has $a\ne b$ and either
    \begin{align}
        a_1=b_1\in\{p_1,q_1\},
    \end{align}
    or
    \begin{align}
        a_2=b_2\in\{p_2,q_2\}.
    \end{align}
    No layer-membership or late-scale assumption is required.
\end{theorem}
\begin{proof}\leanok
    \uses{thm:area_law_unsplit_side_endpoints,
        thm:area_law_elementary_side_geometry}
    Each whole side has length $t>0$ and a constant coordinate equal
    to a boundary coordinate of the square. An elementary segment is
    either this whole side or one of its two midpoint halves. Its
    endpoints retain that constant coordinate and differ in the other
    coordinate by $t$ or $t/2$.
\end{proof}

\begin{theorem}[Unique opposing fine cell or dummy neighborhood]
    \label{thm:area_law_unique_elementary_side_opponent}
    \lean{TNLean.PEPS.AreaLaw.Geometry.exists_unique_elementarySide_opponent}
    \leanok
    \uses{def:area_law_dyadic_layers,def:area_law_fine_layer_indices,
        def:area_law_fine_belt_scales,def:area_law_actual_side_mask}
    Fix a common origin $o\in\mathbb R^2$, finite endpoint set $Z$,
    integer width $C_0\ge2$ and natural indices $k_0\le k$ with
    $k\ge50\,000\,000$. Let $z\in F_{k,\ell_k}$ index an actual fine
    cell. For any elementary segment $e=[a_i,b_i]$ of this cell under
    its actual midpoint subdivision relative to $k_0$, consider the
    identifiers
    \begin{align}
        \mathcal I=\{\star\}\sqcup
        \{(h,w):h\ge k_0,\ w\in F_{h,\ell_h},\ (h,w)\ne(k,z)\}.
    \end{align}
    Associate to them the closed regions
    \begin{align}
        R_\star=\overline{N_{k_0}},\qquad
        R_{(h,w)}=\overline{Q_{\ell_h}(w)}.
    \end{align}
    There is exactly one $P\in\mathcal I$ such that $e\subseteq R_P$.
    The dummy neighborhood is a single region even when several coarse
    cells contribute to it. No further lower bound on the opposing
    layer index is imposed beyond $h\ge k_0$.
\end{theorem}
\begin{proof}\leanok
    \uses{thm:area_law_elementary_side_opponent,
        thm:area_law_elementary_side_boundary_geometry,
        thm:area_law_fine_cells_disjoint,thm:area_law_dummy_layer_disjoint,
        thm:area_law_fine_cell_containment,thm:area_law_dyadic_closures,
        thm:area_law_elementary_corner_endpoint,
        thm:area_law_dummy_corner_elementary_side,
        thm:area_law_cell_fan_mark_vertices,
        thm:area_law_belt_marks_closed_cell}
    Existence follows from the preceding opposing-region theorem.
    To prove uniqueness, put $x=(a_i+b_i)/2$. The segment is
    nondegenerate and lies on one side of the reference square, so
    $x$ is interior to that side in its tangential coordinate.

    Suppose first that two distinct actual fine cells contain the
    whole segment in their closures. Their interiors and the interior
    of the reference cell are pairwise disjoint. Neither opposing
    square can have a corner at $x$, since an opposing corner on an
    elementary segment must be an endpoint. Consequently the
    tangential coordinate of $x$ lies strictly between the two
    tangential bounds of each opposing square.

    Write $x=(\xi,\eta)$ in tangential and normal coordinates, and
    write the closures of the reference square and the two opposing
    squares as $[u_j,v_j]\times[u'_j,v'_j]$ for $j=0,1,2$. Then
    \begin{align}
        \xi  & \in(u_0,v_0)\cap(u_1,v_1)\cap(u_2,v_2),   \\
        \eta & \in[u'_0,v'_0]\cap[u'_1,v'_1]\cap[u'_2,v'_2].
    \end{align}
    Since the tangential open intervals overlap, pairwise disjoint
    interiors force the normal open intervals $(u'_j,v'_j)$ to be
    pairwise disjoint. Each of them has positive length and contains
    $(\eta-\epsilon,\eta)$ or $(\eta,\eta+\epsilon)$ for some
    $\epsilon>0$. Two of the three therefore meet, a contradiction.
    Interchanging the coordinates handles a horizontal elementary
    segment.

    If both the dummy neighborhood and a distinct actual fine cell
    contain the segment, the finite union defining the dummy
    neighborhood supplies a contributing closed coarse cell
    containing $x$. Its interior is disjoint from both fine-cell
    interiors because their layer indices are at least $k_0$.
    A dummy corner on the elementary segment is again an endpoint,
    so $x$ is not such a corner. The same three-rectangle argument
    gives a contradiction. Thus the existing opposing identifier
    is unique.
\end{proof}

These assertions identify the opposing actual fine cell or dummy
neighborhood. Their composition with primary identifiers and fan-run
labels, the isolated-star construction and the complete two-family
partition remain further obligations.
